% ==============================================================================
%  1.5  El modelo lineal. Inferencias sobre una, dos y varias medias.
%       Introducción al ANOVA.
% ==============================================================================
\section{El modelo lineal para inferencias sobre medias. Introducción al ANOVA}
\label{sec:t1-medias}

Esta sección trata los métodos de \textbf{inferencia sobre las medias} de una, dos y varias
poblaciones. Estas inferencias pueden basarse en los principios de las distribuciones
muestrales o en \textbf{modelos lineales}. En esta asignatura interesan los modelos lineales,
que también se apoyan en las \textbf{distribuciones muestrales}, aunque con un enfoque algo
distinto.

\paragraph{Proceso de ``hacer inferencia''.}
\begin{enumerate}
  \item Se identifican uno o más \textbf{parámetros} que caracterizan la población de interés.
  \item Se toma una \textbf{muestra representativa} de la población.
  \item Para estimar el parámetro se elige un \textbf{estadístico} que pueda calcularse con la
    muestra. Al calcularse a partir de una muestra aleatoria, es una variable aleatoria.
  \item Se determina su \textbf{distribución en el muestreo}, es decir, la distribución
    teórica del estadístico en las distintas muestras posibles, que describe el
    comportamiento de todos sus valores.
  \item Con esa distribución y el valor observado en la muestra, se \textbf{hace inferencia}
    sobre el parámetro.
\end{enumerate}

\begin{definicion}[Suma de cuadrados y media cuadrática][def:SS-MS]
  \begin{itemize}
    \item \textbf{Suma de cuadrados} ($SS$): suma de las desviaciones cuadráticas de los valores
      observados respecto de la media muestral,
      \[ SS_{xx} = \sumi (x_i-\media{x})^2 . \]
    \item \textbf{Media cuadrática} ($MS$): suma de cuadrados dividida por sus grados de
      libertad,
      \[ MS_x = \frac{SS_{xx}}{n-1} = \frac{\sumi (x_i-\media{x})^2}{n-1} = S^2 , \]
      es decir, la varianza muestral corregida $S^2$ es una media cuadrática.
  \end{itemize}
\end{definicion}

\begin{nota}
  Análogamente, $SS_{xy}=\sumi(x_i-\media{x})(y_i-\media{y})$. Con esta notación, el
  coeficiente de correlación de la \cref{def:pearson} es
  $r_{xy} = SS_{xy}/\sqrt{SS_{xx}\,SS_{yy}}$.
\end{nota}

Los tres contrastes que se estudian a continuación siguen el mismo esquema:
\begin{resumen}[Esquema común: modelo restringido frente a no restringido]
  \begin{enumerate}
    \item Se plantea un modelo lineal \textbf{no restringido} (compatible con $H_1$) y el modelo
      \textbf{restringido} que impone $H_0$.
    \item Se estima cada uno por \textbf{mínimos cuadrados} y se calcula su suma de cuadrados
      de los errores. Como el modelo restringido es un caso particular del no restringido,
      $\SC{restringido} \ge \SC{no restringido}$.
    \item El aumento $\SC{hipótesis}=\SC{restringido}-\SC{no restringido}$ mide cuánto empeora
      el ajuste al imponer $H_0$. Si $H_0$ es cierta, este aumento debería ser pequeño.
    \item Se comparan $\MC{hipótesis}$ y $\MC{no restringido}$ con un estadístico $F$ y se
      rechaza $H_0$ si $F$ es grande.
  \end{enumerate}
\end{resumen}

% ------------------------------------------------------------------------------
\subsection{Inferencias para una media}
\label{sec:t1-una-media}

Se parte de una muestra aleatoria de tamaño $n$ de una población normal con media $\mu$ y
desviación típica $\sigma$, y el objetivo es hacer inferencias sobre $\mu$. Para ello se
plantea el \textbf{modelo lineal}
\[
  y_i = \mu + \varepsilon_i, \qquad i=1,\ldots,n,
\]
donde $y_i$ es el $i$-ésimo valor observado de la respuesta $Y$, $\mu$ es la media poblacional
y $\varepsilon_1,\ldots,\varepsilon_n$ son variables aleatorias independientes
$\Normal(0,\sigma^2)$.
\begin{itemize}
  \item \textbf{Componente sistemática}: $\mu$. En ausencia de variabilidad ($\sigma=0$),
    $y_i=\mu$ para todo $i$. Como la media de los $\varepsilon_i$ es 0, $\E[Y]=\mu$.
  \item \textbf{Componente aleatoria} (``error''): los $\varepsilon_i=y_i-\mu$, que describen la
    variabilidad de los valores individuales alrededor de $\mu$.
  \item \textbf{Parámetro $\sigma^2$}: es la varianza de $\varepsilon$ y mide la dispersión del
    error, es decir, la variabilidad del modelo. Sirve como \emph{medida del ajuste}: una
    varianza pequeña indica que los valores observados están cerca de la media, y por tanto
    un buen ajuste.
\end{itemize}

\paragraph{Estimación por mínimos cuadrados.} El criterio más utilizado para estimar $\mu$ es
el \textbf{criterio de mínimos cuadrados}, que consiste en buscar el valor $\est{\mu}$ que
minimiza la suma de los cuadrados de las diferencias entre los valores observados y la media
estimada,
\[
  SS_{yy} = \sumi \est{\varepsilon}_i^{\,2} = \sumi (y_i-\est{\mu})^2 .
\]
Derivando respecto de $\mu$ e igualando a cero,
\[
  \frac{\partial}{\partial \mu}\sumi (y_i-\mu)^2 = -2\sumi y_i + 2n\mu = 0
  \quad\Longrightarrow\quad
  \est{\mu} = \frac{\sumi y_i}{n} = \media{y} .
\]
Por tanto, $SS_{yy}=\sumi (y_i-\media{y})^2$, que es el numerador de $s_y^2$, la estimación de
$\sigma^2$. Es decir, $\est{\mu}=\media{y}$ minimiza la variabilidad del modelo y proporciona
el modelo que mejor se ajusta a los datos.

\paragraph{Contraste $H_0\colon \mu=\mu_0$ frente a $H_1\colon \mu\neq\mu_0$.} Hay que elegir
entre dos modelos:
\begin{enumerate}
  \item \textbf{Modelo restringido} (verifica $H_0$): $y_i = \mu_0 + \varepsilon_i$.
  \item \textbf{Modelo no restringido} ($\mu$ libre; apoyaría la alternativa):
    $y_i = \mu + \varepsilon_i$.
\end{enumerate}
Se elige el que minimiza la varianza. Las sumas de cuadrados de ambos modelos son
\[
  \SC{restringido} = \sumi (y_i-\mu_0)^2, \qquad
  \SC{no restringido} = \sumi (y_i-\media{y})^2 .
\]
Como $\media{y}$ es el estimador mínimo cuadrático, minimiza la suma de cuadrados y
$\SC{restringido} \ge \SC{no restringido}$. La distancia entre ambas es la base del contraste:
\begin{align*}
  \SC{restringido} &= \sumi (y_i-\mu_0)^2 = \sumi (y_i-\media{y}+\media{y}-\mu_0)^2 \\
  &= \sumi (y_i-\media{y})^2 + \sumi (\media{y}-\mu_0)^2
     + 2(\media{y}-\mu_0)\underbrace{\sumi (y_i-\media{y})}_{=\,0} \\
  &= \SC{no restringido} + n(\media{y}-\mu_0)^2 .
\end{align*}

\begin{proposicion}[Descomposición de la variabilidad (una media)][prop:desc-una-media]
  \[
    \underbrace{\sumi (y_i-\mu_0)^2}_{\SC{restringido}}
    = \underbrace{\sumi (y_i-\media{y})^2}_{\SC{no restringido}}
    + \underbrace{n(\media{y}-\mu_0)^2}_{\SC{hipótesis}}
  \]
  \begin{itemize}
    \item $\SC{restringido}$ tiene $n$ grados de libertad: es la suma de $n$ desviaciones
      respecto de una cantidad, $\mu_0$, que no se calcula a partir de los datos.
    \item $\SC{no restringido}$ tiene $n-1$ grados de libertad: suma de $n$ desviaciones
      respecto de $\media{y}$, estimada por mínimos cuadrados.
    \item $\SC{hipótesis}$ es lo que aumenta el error al imponer que $H_0$ es cierta; tiene
      1 grado de libertad.
  \end{itemize}
  La igualdad también se cumple para los grados de libertad: $n=(n-1)+1$.
\end{proposicion}

\paragraph{Medias cuadráticas y estadístico $F$.}
\[
  \MC{no restringido} = \frac{\sumi (y_i-\media{y})^2}{n-1}, \qquad
  \MC{hipótesis} = \frac{n(\media{y}-\mu_0)^2}{1}.
\]
Ambas estiman la misma varianza $\sigma^2$ si $H_0$ es cierta, y entonces
\[
  F = \frac{\MC{hipótesis}}{\MC{no restringido}} \bajoHo \Fdist{1}{n-1} .
\]
Si $H_0$ no es cierta, el numerador (y por tanto $F$) tiende a ser mayor. \textbf{Se rechaza
$H_0$ si $F_{\text{obs}} > F_{1,\,n-1;\,1-\alpha}$}.

\begin{nota}
  Como $\MC{no restringido}=s^2$, se tiene $F=\left(\dfrac{\media{y}-\mu_0}{s/\sqrt{n}}\right)^2=t^2$:
  es el cuadrado del estadístico $t$ habitual para una media, y $F_{1,n-1}$ es la distribución
  de $t_{n-1}^2$. Ambos contrastes son equivalentes.
\end{nota}

% ------------------------------------------------------------------------------
\subsection{Inferencias para dos medias}
\label{sec:t1-dos-medias}

Se consideran dos poblaciones de una variable $Y$ normal, con medias $\mu_1$, $\mu_2$ y
varianzas $\sigma_1^2$, $\sigma_2^2$, de las que se obtienen muestras aleatorias de tamaños
$n_1$ y $n_2$, con valores observados $y_{ij}$, $i=1,2$, $j=1,\ldots,n_i$. El interés está en
la diferencia de medias $\delta=\mu_1-\mu_2$. El \textbf{modelo lineal} es
\[
  y_{ij} = \mu_i + \varepsilon_{ij},
\]
donde $y_{ij}$ es el $j$-ésimo valor observado en la población $i$, $\mu_i$ es la media de la
población $i$ y los $\varepsilon_{ij}$ son observaciones de una variable aleatoria
$\Normal(0,\sigma^2)$. Obsérvese que \textbf{se supone que las dos poblaciones tienen la misma
varianza} $\sigma^2$ ($\sigma_1^2=\sigma_2^2=\sigma^2$).

\paragraph{Contraste $H_0\colon \mu_1=\mu_2$ frente a $H_1\colon \mu_1\neq\mu_2$.}
\begin{itemize}
  \item \textbf{No restringido}: $y_{ij}=\mu_i+\varepsilon_{ij}$. Minimizando
    $\sum_i\sum_j (y_{ij}-\est{\mu}_i)^2$ se obtienen $\est{\mu}_1=\media{y}_1$ y
    $\est{\mu}_2=\media{y}_2$ (las medias de cada muestra).
  \item \textbf{Restringido} ($H_0$ cierta): $y_{ij}=\mu+\varepsilon_{ij}$, con estimador
    mínimo cuadrático $\est{\mu}=\media{y}$ (la media de todas las observaciones).
\end{itemize}

\begin{proposicion}[Descomposición de la variabilidad (dos medias)][prop:desc-dos-medias]
  $\SC{restringido} = \SC{no restringido} + \SC{hipótesis}$, donde
  \begin{itemize}
    \item $\SC{restringido} = \sum_i\sum_j (y_{ij}-\media{y})^2$, con $n_1+n_2-1$ g.l.;
    \item $\SC{no restringido} = \sum_j (y_{1j}-\media{y}_1)^2 + \sum_j (y_{2j}-\media{y}_2)^2
      = SS_1+SS_2$, con $n_1+n_2-2$ g.l.;
    \item $\SC{hipótesis} = n_1(\media{y}_1-\media{y})^2 + n_2(\media{y}_2-\media{y})^2$, con
      1 g.l. (modelo sin restricciones: 2 parámetros; con restricciones: 1 parámetro).
  \end{itemize}
\end{proposicion}

\paragraph{Medias cuadráticas y estadístico $F$.}
\[
  \MC{no restringido} = \frac{SS_1+SS_2}{n_1+n_2-2}, \qquad
  \MC{hipótesis} = n_1(\media{y}_1-\media{y})^2 + n_2(\media{y}_2-\media{y})^2 .
\]
Ambas estiman $\sigma^2$ si $H_0$ es cierta, y entonces
\[
  F = \frac{\MC{hipótesis}}{\MC{no restringido}} \bajoHo \Fdist{1}{n_1+n_2-2} .
\]
Si $H_0$ no es cierta, $F$ tiende a ser mayor. \textbf{Se rechaza $H_0$ si
$F_{\text{obs}} > F_{1,\,n_1+n_2-2;\,1-\alpha}$}.

\begin{nota}
  $\MC{no restringido}$ es la varianza combinada $s_p^2$ del contraste $t$ para dos muestras
  con varianzas iguales, y $F=t^2$. De nuevo, el enfoque del modelo lineal conduce al contraste
  clásico.
\end{nota}

% ------------------------------------------------------------------------------
\subsection{Inferencias para varias medias: el ANOVA}
\label{sec:t1-anova}

El modelo anterior se generaliza al caso de $k>2$ medias. Suponiendo que los datos son
muestras aleatorias independientes de tamaños $n_i$ de $k$ poblaciones normales, el modelo
lineal es
\[
  y_{ij} = \mu_i + \varepsilon_{ij}, \qquad i=1,\ldots,k,\quad j=1,\ldots,n_i,
\]
donde $y_{ij}$ es el $j$-ésimo valor observado en la población $i$, $\mu_i$ es la media de la
población $i$ y los $\varepsilon_{ij}$ son observaciones de una variable aleatoria
$\Normal(0,\sigma^2)$ (misma varianza en todas las poblaciones).

\begin{definicion}[ANOVA][def:anova]
  Este modelo se conoce como \textbf{análisis de la varianza (ANOVA)}. Se denota por
  $N=\sum_{i=1}^k n_i$ el número total de observaciones.
\end{definicion}

\paragraph{Contraste.}
\[
  H_0\colon \mu_i=\mu_j \ \ \forall i\neq j
  \qquad\text{frente a}\qquad
  H_1\colon \mu_i\neq\mu_j \ \text{para uno o más pares}.
\]
\begin{itemize}
  \item \textbf{No restringido}: minimizando $\sum_i\sum_j (y_{ij}-\est{\mu}_i)^2$ se obtiene
    $\est{\mu}_i=\media{y}_i$, $i=1,\ldots,k$.
  \item \textbf{Restringido}: $y_{ij}=\mu+\varepsilon_{ij}$, con $\est{\mu}=\media{y}$.
\end{itemize}

\begin{proposicion}[Descomposición de la variabilidad (ANOVA)][prop:desc-anova]
  $\SC{restringido} = \SC{no restringido} + \SC{hipótesis}$, donde
  \begin{itemize}
    \item $\SC{restringido} = \sum_i\sum_j (y_{ij}-\media{y})^2$, con $N-1$ g.l.;
    \item $\SC{no restringido} = \sum_j (y_{1j}-\media{y}_1)^2+\cdots+\sum_j (y_{kj}-\media{y}_k)^2
      = SS_1+\cdots+SS_k$, con $N-k$ g.l.;
    \item $\SC{hipótesis} = \sum_i n_i(\media{y}_i-\media{y})^2$, con $k-1$ g.l. (modelo sin
      restricciones: $k$ parámetros; con restricciones: 1 parámetro).
  \end{itemize}
\end{proposicion}

\paragraph{Medias cuadráticas y estadístico $F$.}
\[
  \MC{no restringido} = \frac{SS_1+\cdots+SS_k}{N-k}, \qquad
  \MC{hipótesis} = \frac{\sum_i n_i(\media{y}_i-\media{y})^2}{k-1}.
\]
Ambas estiman $\sigma^2$ si $H_0$ es cierta, y entonces
\[
  F = \frac{\MC{hipótesis}}{\MC{no restringido}} \bajoHo \Fdist{k-1}{N-k} .
\]
Si $H_0$ no es cierta, $F$ tiende a ser mayor. \textbf{Se rechaza $H_0$ si
$F_{\text{obs}} > F_{k-1,\,N-k;\,1-\alpha}$}.

\begin{nota}
  El caso de dos medias es el ANOVA con $k=2$ y $N=n_1+n_2$.
\end{nota}

\begin{resumen}[Resumen de los tres contrastes]
  \centering
  \small
  \renewcommand{\arraystretch}{1.35}
  \begin{tabular}{lccc}
    \toprule
    & \textbf{Una media} & \textbf{Dos medias} & \textbf{$k$ medias (ANOVA)} \\
    \midrule
    $H_0$ & $\mu=\mu_0$ & $\mu_1=\mu_2$ & $\mu_1=\cdots=\mu_k$ \\
    \midrule
    $\SC{restringido}$ & $\sum_i(y_i-\mu_0)^2$ & $\sum_i\sum_j(y_{ij}-\media{y})^2$
      & $\sum_i\sum_j(y_{ij}-\media{y})^2$ \\
    \quad g.l. & $n$ & $n_1+n_2-1$ & $N-1$ \\
    $\SC{no restringido}$ & $\sum_i(y_i-\media{y})^2$ & $SS_1+SS_2$ & $SS_1+\cdots+SS_k$ \\
    \quad g.l. & $n-1$ & $n_1+n_2-2$ & $N-k$ \\
    $\SC{hipótesis}$ & $n(\media{y}-\mu_0)^2$ & $\sum_{i=1}^{2} n_i(\media{y}_i-\media{y})^2$
      & $\sum_{i=1}^{k} n_i(\media{y}_i-\media{y})^2$ \\
    \quad g.l. & $1$ & $1$ & $k-1$ \\
    \midrule
    $F$ bajo $H_0$ & $\Fdist{1}{n-1}$ & $\Fdist{1}{n_1+n_2-2}$ & $\Fdist{k-1}{N-k}$ \\
    \bottomrule
  \end{tabular}
  \par\medskip\raggedright
  En todos los casos $F=\MC{hipótesis}/\MC{no restringido}$ y se rechaza $H_0$ si
  $F_{\text{obs}}$ supera el cuantil $1-\alpha$ de la distribución $F$ correspondiente.
\end{resumen}
